% =====================================================================
\chapter{Soluciones de los ejercicios}
\label{ap:soluciones}
% =====================================================================

Cada solución se presenta en cinco partes: datos, planteamiento, desarrollo, resultado y comprobación. En los pasos intermedios se redondea a cuatro cifras significativas salvo que se indique otra cosa; los resultados exactos se dan como fracción cuando existe.

% ---------------------------------------------------------------------
\section{Tema 1}

\begin{solucion}{eje:1-tres}
\datos tres formas de adaptar un modelo: pocos ejemplos, RAG y LoRA; un documento publicado ayer.
\planteamiento clasificar cada técnica según modifique el contexto o los parámetros y según quién elija el contenido añadido.
\pasos
\begin{enumerate}[label=\arabic*.]
  \item Pocos ejemplos: se añaden $k$ demostraciones al \emph{prompt}. Los parámetros $\theta$ no cambian; solo cambia la secuencia de entrada (\cref{def:icl}).
  \item RAG: también cambia solo el contexto, pero el contenido lo elige un recuperador sobre una colección, que se actualiza indexando documentos, sin entrenar.
  \item LoRA: se entrenan las matrices $B$ y $A$ (módulo~20); cambian unos pocos parámetros efectivos y el contexto no.
  \item Documento de ayer: RAG lo usa en cuanto se indexa; con pocos ejemplos también podría usarse si alguien lo pega a mano en el \emph{prompt}; LoRA necesitaría un entrenamiento nuevo.
\end{enumerate}
\resultado \textbf{pocos ejemplos y RAG modifican el contexto; LoRA modifica parámetros. RAG aprovecha el documento de ayer sin entrenar.}
\comprobacion si se borra el \emph{prompt}, los efectos de los pocos ejemplos y de RAG desaparecen y el de LoRA permanece, porque está en los pesos.
\end{solucion}

\begin{solucion}{eje:1-induccion}
\datos secuencia $x_{1:6}=(\texttt{a},\texttt{b},\texttt{c},\texttt{a},\texttt{d},\texttt{a})$; circuito de la \cref{prop:induccion}.
\planteamiento hallar $J=\{j\in\{2,\dots,t\}:x_{j-1}=x_t\}$ y promediar los $e(x_j)$ con $j\in J$.
\pasos
\begin{enumerate}[label=\arabic*.]
  \item (a) $x_6=\texttt{a}$. Las posiciones con $x_{j-1}=\texttt{a}$ son $j=2$ ($x_1=\texttt{a}$) y $j=5$ ($x_4=\texttt{a}$); $j=7$ no existe. Luego $J=\{2,5\}$.
  \item $o_6=\tfrac12\bigl(e(x_2)+e(x_5)\bigr)=\tfrac12\bigl(e(\texttt{b})+e(\texttt{d})\bigr)$: probabilidad $1/2$ para \texttt{b} y para \texttt{d}.
  \item (b) Con $x_6=\texttt{c}$, la única posición con $x_{j-1}=\texttt{c}$ es $j=4$ ($x_3=\texttt{c}$), luego $J=\{4\}$ y $o_6=e(x_4)=e(\texttt{a})$.
  \item (c) Si $J=\varnothing$, todas las puntuaciones valen~0, la softmax es uniforme sobre las posiciones y la salida es la media de todos los valores: la cabeza no copia nada concreto.
\end{enumerate}
\resultado \textbf{(a) $J=\{2,5\}$ y $o_6=\frac12(e(\texttt{b})+e(\texttt{d}))$; (b) $J=\{4\}$ y $o_6=e(\texttt{a})$; (c) salida uniforme, sin copia.}
\comprobacion en el contexto, \texttt{a} va seguida una vez de \texttt{b} y otra de \texttt{d}, y \texttt{c} va seguida de \texttt{a}: son las frecuencias de un modelo de bigramas estimado sobre la propia secuencia.
\end{solucion}

\begin{solucion}{eje:1-coste}
\datos $N=7\cdot 10^{9}$; 1000 consultas; 200 tokens de razonamiento; $m=10$; 5 demostraciones de 40 tokens; $2N$ operaciones por token.
\planteamiento multiplicar el número de tokens por el coste por token.
\pasos
\begin{enumerate}[label=\arabic*.]
  \item Coste por token: $2N=1{,}4\cdot 10^{10}$ operaciones.
  \item (a) Por consulta: $200\cdot 1{,}4\cdot 10^{10}=2{,}8\cdot 10^{12}$. Para 1000 consultas: $2{,}8\cdot 10^{15}$.
  \item (b) Con $m=10$ razonamientos: $10\cdot 2{,}8\cdot 10^{12}=2{,}8\cdot 10^{13}$ por consulta y $2{,}8\cdot 10^{16}$ en total.
  \item (c) Tokens de demostraciones: $5\cdot 40=200$ por consulta y $200\cdot 1000=2\cdot 10^{5}$ en total. Operaciones: $2\cdot 10^{5}\cdot 1{,}4\cdot 10^{10}=2{,}8\cdot 10^{15}$.
\end{enumerate}
\resultado \textbf{(a) $2{,}8\cdot 10^{15}$ operaciones; (b) $2{,}8\cdot 10^{16}$; (c) $2\cdot 10^{5}$ tokens y $2{,}8\cdot 10^{15}$ operaciones.}
\comprobacion (a) y (c) coinciden porque ambos casos suponen 200 tokens por consulta; la autoconsistencia multiplica el coste por $m$.
\end{solucion}

\begin{solucion}{eje:1-voto}
\datos $p=0{,}6$; votos independientes; $m$ impar.
\planteamiento el número de aciertos es binomial $(m,p)$; el voto acierta si hay más de $m/2$ aciertos. Para (b), imponer $e^{-2m(p-1/2)^{2}}\le 0{,}05$.
\pasos
\begin{enumerate}[label=\arabic*.]
  \item $m=3$: $p^{3}+3p^{2}(1-p)=0{,}216+3\cdot 0{,}36\cdot 0{,}4=0{,}216+0{,}432=0{,}648$.
  \item $m=5$: $p^{5}+5p^{4}(1-p)+10p^{3}(1-p)^{2}=0{,}07776+5\cdot 0{,}1296\cdot 0{,}4+10\cdot 0{,}216\cdot 0{,}16=0{,}07776+0{,}2592+0{,}3456=0{,}68256$.
  \item (b) $(p-1/2)^{2}=0{,}01$. La condición es $e^{-0{,}02m}\le 0{,}05$, es decir, $0{,}02m\ge\ln 20=2{,}9957$, o $m\ge 149{,}79$.
  \item El menor impar que cumple la condición es $m=151$.
\end{enumerate}
\resultado \textbf{(a) $0{,}648$ con $m=3$ y $0{,}6826$ con $m=5$; (b) $m=151$ garantiza un error de a lo sumo $0{,}05$.}
\comprobacion el cálculo binomial exacto muestra que basta $m=67$ (error $0{,}0485$) y que con $m=151$ el error es $0{,}0064$: la cota de Hoeffding es válida pero conservadora.
\end{solucion}

\begin{solucion}{eje:1-kv}
\datos 40 capas; $d=5120$; 2 bytes por número; 32\,000 tokens; con GQA, 8 cabezas de clave y valor de 40.
\planteamiento por la \cref{prop:kv}, bytes por token $=2\,n_{\mathrm{capas}}\,d_{kv}\,\beta$, con $d_{kv}=d$ sin GQA.
\pasos
\begin{enumerate}[label=\arabic*.]
  \item $2\cdot 40=80$; $80\cdot 5120=409\,600$; $409\,600\cdot 2=819\,200$ bytes por token, es decir, $819\,200/1\,048\,576=0{,}7813$~MiB.
  \item Con 32\,000 tokens: $819\,200\cdot 32\,000=2{,}621\cdot 10^{10}$ bytes $=26{,}21$~GB.
  \item GQA: $d_{\mathrm{cabeza}}=5120/40=128$ y $d_{kv}=8\cdot 128=1024$, cinco veces menos.
  \item Bytes por token con GQA: $2\cdot 40\cdot 1024\cdot 2=163\,840$; con 32\,000 tokens, $5{,}243\cdot 10^{9}$ bytes $=5{,}24$~GB.
\end{enumerate}
\resultado \textbf{819\,200 bytes por token ($0{,}78$~MiB) y $26{,}2$~GB con 32\,000 tokens; con GQA, $5{,}24$~GB.}
\comprobacion $26{,}21/5=5{,}24$, como corresponde a dividir $d_{kv}$ por~5.
\end{solucion}

\begin{solucion}{eje:1-gd}
\datos $L(w)=\frac12\sum_i(w^{\top}x_i-y_i)^{2}$, $w_0=0$, tasa $\eta$; $x_1=(1,0)$, $y_1=2$, $x_2=(0,1)$, $y_2=-1$, $x_q=(1,1)$, $\eta=0{,}5$.
\planteamiento calcular $w_1$ y comparar $w_1^{\top}x_q$ con la atención lineal $\sum_i(q^{\top}k_i)v_i$.
\pasos
\begin{enumerate}[label=\arabic*.]
  \item Por la regla de la cadena, $\nabla_w\frac12(w^{\top}x_i-y_i)^{2}=(w^{\top}x_i-y_i)\,x_i$; sumando, $\nabla L(w)=\sum_i(w^{\top}x_i-y_i)\,x_i$.
  \item En $w_0=0$: $\nabla L(0)=-\sum_i y_ix_i$.
  \item Paso: $w_1=0-\eta\nabla L(0)=\eta\sum_i y_ix_i$.
  \item Predicción: $w_1^{\top}x_q=\eta\sum_i y_i\,(x_i^{\top}x_q)$, por linealidad.
  \item Atención lineal con $q=x_q$, $k_i=x_i$, $v_i=y_i$: $\sum_i(x_q^{\top}x_i)\,y_i$; multiplicada por $\eta$ coincide con el paso~4.
  \item Ejemplo: $\sum_i y_ix_i=2(1,0)-(0,1)=(2,-1)$; $w_1=(1;\,-0{,}5)$; $w_1^{\top}x_q=1-0{,}5=0{,}5$.
  \item Por la atención: $x_1^{\top}x_q=1$, $x_2^{\top}x_q=1$; $0{,}5\,(2\cdot 1+(-1)\cdot 1)=0{,}5$.
  \item Elementos ausentes: la softmax (aquí los pesos $q^{\top}k$ no se normalizan), las proyecciones aprendidas $W_Q$, $W_K$, $W_V$ y de salida, las varias cabezas, los bloques MLP, la normalización, el apilamiento de capas y la máscara causal.
\end{enumerate}
\resultado \textbf{$\hat y_q=0{,}5$ por las dos vías.}
\comprobacion $w_1^{\top}x_1=1$ (objetivo~2) y $w_1^{\top}x_2=-0{,}5$ (objetivo $-1$): con $x_i$ ortonormales, cada predicción recorre la mitad del camino desde~0, como corresponde a $\eta=0{,}5$.
\end{solucion}

\begin{solucion}{eje:1-esfera}
\datos $\norm{q}=\norm{k_i}=1$; $\kappa_\sigma(q,k)=\exp(-\norm{q-k}^{2}/(2\sigma^{2}))$; $s_i=q\cdot k_i$.
\planteamiento desarrollar el cuadrado de la norma y estudiar los pesos $w_i=\softmax(s/\sigma^{2})_i$ en los dos límites.
\pasos
\begin{enumerate}[label=\arabic*.]
  \item (a) $\norm{q-k_i}^{2}=\norm{q}^{2}-2q\cdot k_i+\norm{k_i}^{2}=2-2s_i$. Entonces $\kappa_\sigma(q,k_i)=e^{-1/\sigma^{2}}e^{s_i/\sigma^{2}}$, y el factor común se cancela en $\kappa_\sigma(q,k_i)/\sum_l\kappa_\sigma(q,k_l)$: queda $\softmax(s/\sigma^{2})_i$.
  \item (b) Sea $s^\ast=s_{i^\ast}=\max_i s_i$. Dividiendo numerador y denominador por $e^{s^\ast/\sigma^{2}}$, $w_i=e^{(s_i-s^\ast)/\sigma^{2}}/\sum_l e^{(s_l-s^\ast)/\sigma^{2}}$.
  \item Para $i\ne i^\ast$, $s_i-s^\ast<0$ y $e^{(s_i-s^\ast)/\sigma^{2}}\to 0$ cuando $\sigma\to 0$; para $i^\ast$, el término vale~1. Luego $w_{i^\ast}\to 1$, $w_i\to 0$ para el resto y $\hat f(q)\to v_{i^\ast}$.
  \item Cuando $\sigma\to\infty$, $s_i/\sigma^{2}\to 0$ para todo $i$, cada exponencial tiende a~1 y $w_i\to 1/n$: $\hat f(q)\to\frac1n\sum_i v_i$.
  \item (c) Con demostraciones como claves y valores: con $\sigma\to 0$, el modelo copia la etiqueta de la demostración más parecida a la entrada (vecino más próximo, como la copia de una cabeza de inducción); con $\sigma\to\infty$, ignora la entrada y devuelve la etiqueta media de las demostraciones.
\end{enumerate}
\resultado \textbf{$\hat f(q)=\sum_i\softmax(q\cdot k/\sigma^{2})_iv_i$; tiende a $v_{i^\ast}$ si $\sigma\to 0$ y a la media de los valores si $\sigma\to\infty$.}
\comprobacion con los datos del \cref{ej:nw}, $i^\ast=3$ ($s_3=0{,}96$): $\hat f\to 1$ si $\sigma\to 0$; $\hat f=0{,}7870$ con $\sigma=1$; $\hat f\to(2-1+1)/3=0{,}6667$ si $\sigma\to\infty$. Ambos valores extremos acotan al intermedio.
\end{solucion}

% ---------------------------------------------------------------------
\section{Tema 2}

\begin{solucion}{eje:2-idf}
\datos $\idf_t=\ln(N/\df_t)$; peso RSJ $\ln\bigl((N-\df_t+0{,}5)/(\df_t+0{,}5)\bigr)$ y su variante con $1+$.
\planteamiento estudiar el signo de cada logaritmo y aplicar la \cref{prop:escala}.
\pasos
\begin{enumerate}[label=\arabic*.]
  \item (a) Si $\df_t=N$, $\idf_t=\ln 1=0$ y el peso de $t$ es~0 en la consulta y en todos los documentos: no contribuye al producto escalar ni a las normas, como si no existiera. No ayuda a distinguir documentos.
  \item (b) El peso RSJ es negativo si y solo si $N-\df_t+0{,}5<\df_t+0{,}5$, es decir, $\df_t>N/2$. La variante es $\ln(1+x)$ con $x>0$, que es positiva. Con $N=10$ y $\df_t=8$: $-1{,}2238$ frente a $0{,}2578$.
  \item (c) $\log_{10}x=\ln x/\ln 10$: todos los $\idf_t$ quedan multiplicados por $1/\ln 10$. Con pesos $f_{t,d}\cdot\idf_t$, todos los vectores se multiplican por esa constante y, por la \cref{prop:escala}, el coseno no cambia.
  \item Matiz: si además se usa el peso sublineal $1+\log f_{t,d}$, cambiar la base de ese logaritmo no es un cambio de escala uniforme, y el ranking podría cambiar.
\end{enumerate}
\resultado \textbf{(a) idf nula y contribución nula; (b) negativo si $\df_t>N/2$, la variante es siempre positiva; (c) con tf bruta, el ranking no cambia.}
\comprobacion $\ln(1+2{,}5/8{,}5)=\ln 1{,}2941=0{,}2578>0$.
\end{solucion}

\begin{solucion}{eje:2-tfidf}
\datos $D_1=$ «atención query key», $D_2=$ «atención softmax», $D_3=$ «query sql base»; consulta «query atención»; $\idf=\ln(N/\df)$.
\planteamiento construir los vectores tf-idf en el orden (atención, base, key, query, softmax, sql) y calcular cosenos, con los idf redondeados a cuatro decimales.
\pasos
\begin{enumerate}[label=\arabic*.]
  \item $N=3$; $\df(\text{atención})=\df(\text{query})=2$; key, softmax, sql y base tienen $\df=1$.
  \item $\idf(\text{atención})=\idf(\text{query})=\ln 1{,}5=0{,}4055$ y los demás, $\ln 3=1{,}0986$. Cuadrados: $0{,}1644$ y $1{,}2069$.
  \item $V(q)=(0{,}4055;\,0;\,0;\,0{,}4055;\,0;\,0)$, $\norm{V(q)}=\sqrt{0{,}3288}=0{,}5734$.
  \item $V(D_1)=(0{,}4055;\,0;\,1{,}0986;\,0{,}4055;\,0;\,0)$, $\norm{V(D_1)}=\sqrt{1{,}5357}=1{,}2393$; $V(q)\cdot V(D_1)=0{,}3288$; $\cos=0{,}3288/(0{,}5734\cdot 1{,}2393)=0{,}4627$.
  \item $V(D_2)=(0{,}4055;\,0;\,0;\,0;\,1{,}0986;\,0)$, $\norm{V(D_2)}=\sqrt{1{,}3713}=1{,}1710$; producto $0{,}1644$; $\cos=0{,}1644/0{,}6715=0{,}2448$.
  \item $V(D_3)=(0;\,1{,}0986;\,0;\,0{,}4055;\,0;\,1{,}0986)$, $\norm{V(D_3)}=\sqrt{2{,}5783}=1{,}6057$; producto $0{,}1644$; $\cos=0{,}1644/0{,}9207=0{,}1786$.
\end{enumerate}
\resultado \textbf{$\cos=0{,}4627$ ($D_1$), $0{,}2448$ ($D_2$) y $0{,}1786$ ($D_3$); orden $D_1$, $D_2$, $D_3$.}
\comprobacion $D_2$ y $D_3$ comparten un término con la consulta y tienen el mismo producto escalar; $D_3$ queda detrás porque su norma es mayor, al tener más peso en términos ajenos a la consulta.
\end{solucion}

\begin{solucion}{eje:2-iir-car}
\datos $\idf=(1{,}65;\,2{,}08;\,1{,}62;\,1{,}50)$ para (car, auto, insurance, best); frecuencias del enunciado.
\planteamiento $w=f\cdot\idf$; dividir cada columna por su norma; con consulta binaria, la puntuación es la suma de los pesos normalizados de car e insurance.
\pasos
\begin{enumerate}[label=\arabic*.]
  \item (a) Filas car, auto, insurance, best; columnas Doc1, Doc2, Doc3:
  \[
  W=\begin{pmatrix}44{,}55&6{,}60&39{,}60\\6{,}24&68{,}64&0\\0&53{,}46&46{,}98\\21{,}00&0&25{,}50\end{pmatrix}.
  \]
  \item (b) Norma de Doc1: $\sqrt{1984{,}7+38{,}94+441}=\sqrt{2464{,}6}=49{,}65$.
  \item Norma de Doc2: $\sqrt{43{,}56+4711{,}4+2858{,}0}=\sqrt{7613{,}0}=87{,}25$.
  \item Norma de Doc3: $\sqrt{1568{,}2+2207{,}1+650{,}25}=\sqrt{4425{,}5}=66{,}52$.
  \item Normalizando:
  \[
  \hat W=\begin{pmatrix}0{,}8974&0{,}0756&0{,}5953\\0{,}1257&0{,}7867&0\\0&0{,}6127&0{,}7062\\0{,}4230&0&0{,}3833\end{pmatrix}.
  \]
  \item (c) Doc1: $0{,}8974+0=0{,}8974$; Doc2: $0{,}0756+0{,}6127=0{,}6883$; Doc3: $0{,}5953+0{,}7062=1{,}3015$.
\end{enumerate}
\resultado \textbf{orden Doc3 ($1{,}302$), Doc1 ($0{,}897$), Doc2 ($0{,}688$).}
\comprobacion la columna de Doc3 tiene norma~1: $0{,}5953^{2}+0{,}7062^{2}+0{,}3833^{2}=0{,}3544+0{,}4987+0{,}1469=1{,}000$.
\end{solucion}

\begin{solucion}{eje:2-iir-cameras}
\datos $N=10^{7}$; consulta: digital~1, cameras~1; documento: digital~1, video~1, cameras~2; $\df$: $10^{4}$, $10^{5}$, $5\cdot 10^{4}$.
\planteamiento puntuación $=\sum_t q_td_t$ con $q_t=\mathrm{wf}\cdot\idf$ sin normalizar y $d_t=\mathrm{wf}/\norm{\mathrm{wf}}$ sin idf.
\pasos
\begin{enumerate}[label=\arabic*.]
  \item idf: digital $\log_{10}10^{3}=3$; video $\log_{10}10^{2}=2$; cameras $\log_{10}200=2{,}301$.
  \item Consulta en el orden (digital, video, cameras): $q=(3;\,0;\,2{,}301)$.
  \item Documento: $\mathrm{wf}=(1;\,1;\,1+\log_{10}2)=(1;\,1;\,1{,}301)$.
  \item Norma: $\sqrt{1+1+1{,}6927}=\sqrt{3{,}6927}=1{,}9216$; $d=(0{,}5204;\,0{,}5204;\,0{,}6770)$.
  \item Productos: $3\cdot 0{,}5204=1{,}5612$; $0$; $2{,}301\cdot 0{,}6770=1{,}5578$. Suma: $3{,}119$.
\end{enumerate}
\resultado \textbf{puntuación $\approx 3{,}12$.}
\comprobacion si el documento solo contuviera los términos de la consulta, en la dirección de $q$, la puntuación sería $\norm{q}=\sqrt{9+5{,}295}=3{,}781$; sale menos porque «video» se lleva parte de la norma.
\end{solucion}

\begin{solucion}{eje:2-bm25}
\datos $N=10$, $n_t=3$, $k_1=1{,}5$, $b=0{,}75$, $\bar\ell=4$; casos $(f,\abs d)=(1,2),(2,4),(1,8)$.
\planteamiento contribución $=\idf^{\mathrm{BM}}\cdot g(f,\ell)$ con $\ell=\abs d/4$.
\pasos
\begin{enumerate}[label=\arabic*.]
  \item $(N-n_t+0{,}5)/(n_t+0{,}5)=7{,}5/3{,}5=2{,}1429$; $\idf^{\mathrm{BM}}=\ln 3{,}1429=1{,}1451$.
  \item (a) $\ell=0{,}5$; $1-b+b\ell=0{,}25+0{,}375=0{,}625$; denominador $1+1{,}5\cdot 0{,}625=1{,}9375$; $g=2{,}5/1{,}9375=1{,}2903$; contribución $1{,}1451\cdot 1{,}2903=1{,}4776$.
  \item (b) $\ell=1$; denominador $2+1{,}5=3{,}5$; $g=5/3{,}5=1{,}4286$; contribución $1{,}6359$.
  \item (c) $\ell=2$; $0{,}25+1{,}5=1{,}75$; denominador $1+2{,}625=3{,}625$; $g=2{,}5/3{,}625=0{,}6897$; contribución $0{,}7897$.
\end{enumerate}
\resultado \textbf{$\idf^{\mathrm{BM}}=1{,}1451$; (a) $1{,}4776$; (b) $1{,}6359$; (c) $0{,}7897$.}
\comprobacion con $f=1$ y longitud media, $g=1$ y la contribución sería la idf, $1{,}1451$: el trozo corto (a) queda por encima y el largo (c) por debajo. Duplicar $f$ solo multiplica por $1{,}43$.
\end{solucion}

\begin{solucion}{eje:2-sat}
\datos $g(f,\ell)=f(k_1+1)/(f+k_1(1-b+b\ell))$, $k_1>0$, $0\le b\le 1$, $f\ge 0$, $\ell>0$.
\planteamiento escribir $c=k_1(1-b+b\ell)$ y estudiar $g=(k_1+1)f/(f+c)$.
\pasos
\begin{enumerate}[label=\arabic*.]
  \item $c>0$: $1-b\ge 0$ y $b\ell\ge 0$ no se anulan a la vez (si $b=1$, $b\ell=\ell>0$), luego $1-b+b\ell>0$ y $c>0$.
  \item $\partial g/\partial f=(k_1+1)\,[(f+c)-f]/(f+c)^{2}=(k_1+1)\,c/(f+c)^{2}>0$: estrictamente creciente en $f$.
  \item $f/(f+c)=1/(1+c/f)\to 1$ cuando $f\to\infty$; luego $g\to k_1+1$.
  \item $\partial c/\partial\ell=k_1b>0$ si $b>0$.
  \item $\partial g/\partial c=-(k_1+1)f/(f+c)^{2}<0$ si $f>0$. Por la regla de la cadena, $\partial g/\partial\ell=(\partial g/\partial c)(\partial c/\partial\ell)<0$.
\end{enumerate}
\resultado \textbf{$g$ crece con $f$, se satura en $k_1+1$ y decrece con $\ell$ si $b>0$.}
\comprobacion con $k_1=1{,}5$ y $\ell=1$: $g=1;\,1{,}43;\,1{,}67;\,2{,}17$ para $f=1,2,3,10$, creciente y por debajo de $2{,}5$ (\cref{ej:saturacion}).
\end{solucion}

\begin{solucion}{eje:2-rsj}
\datos $c_t=\ln\bigl((N-\df_t)/\df_t\bigr)$, $\idf_t=\ln(N/\df_t)$, $1\le\df_t<N$.
\planteamiento comparar los argumentos de los logaritmos.
\pasos
\begin{enumerate}[label=\arabic*.]
  \item (a) $c_t>0\iff(N-\df_t)/\df_t>1\iff N>2\df_t\iff\df_t<N/2$.
  \item (b) $\idf_t-c_t=\ln\bigl(N/(N-\df_t)\bigr)=-\ln(1-\df_t/N)$. Como $0<1-\df_t/N<1$, el logaritmo es negativo y la diferencia, positiva.
  \item Cuando $\df_t/N\to 0$, $1-\df_t/N\to 1$ y, por continuidad del logaritmo, $-\ln(1-\df_t/N)\to 0$.
  \item (c) $N=10$, $\df_t=3$: $c_t=\ln(7/3)=0{,}8473$ e $\idf_t=\ln(10/3)=1{,}2040$.
  \item $N=10^{6}$, $\df_t=10^{4}$: $c_t=\ln 99=4{,}5951$ e $\idf_t=\ln 100=4{,}6052$.
\end{enumerate}
\resultado \textbf{(a) $c_t>0\iff\df_t<N/2$; (b) $\idf_t-c_t=-\ln(1-\df_t/N)>0$, que tiende a~0; (c) $0{,}8473$ frente a $1{,}2040$, y $4{,}5951$ frente a $4{,}6052$.}
\comprobacion las diferencias son $0{,}3567=-\ln 0{,}7$ y $0{,}0101=-\ln 0{,}99$, como predice~(b).
\end{solucion}

% ---------------------------------------------------------------------
\section{Tema 3}

\begin{solucion}{eje:3-vf}
\planteamiento contrastar cada afirmación con las definiciones y proposiciones del tema.
\pasos
\begin{enumerate}[label=(\alph*)]
  \item \textbf{Verdadera}, por la \cref{prop:monotonia}: $T_k\subseteq T_{k+1}$.
  \item \textbf{Falsa}. Con (R, N): $\Pk{1}=1$ y $\Pk{2}=0{,}5$.
  \item \textbf{Falsa}. Por la \cref{prop:f1}, $F_1\le(P+R)/2$.
  \item \textbf{Verdadera}. $\RR=1/k_1$ depende solo de la posición del primer relevante.
\end{enumerate}
\resultado \textbf{(a) V; (b) F; (c) F; (d) V.}
\comprobacion en el \cref{ej:ranking}, $\Rk{k}=\frac13,\frac13,\frac23,\frac23,\frac23$ no baja y $\Pk{k}=1;\,0{,}5;\,0{,}667;\,0{,}5;\,0{,}4$ sí.
\end{solucion}

\begin{solucion}{eje:3-iir-f1}
\datos devueltos 30; relevantes devueltos 12; relevantes 40; colección 10\,000.
\planteamiento aplicar la \cref{def:pr}.
\pasos
\begin{enumerate}[label=\arabic*.]
  \item $P=12/30=0{,}4$; $R=12/40=0{,}3$.
  \item $F_1=2\cdot 0{,}4\cdot 0{,}3/0{,}7=0{,}24/0{,}7=0{,}3429$.
  \item Sistema que devuelve todo: $R=40/40=1$ y $P=40/10\,000=0{,}004$.
  \item Media aritmética: $(0{,}004+1)/2=0{,}502$. Media armónica: $2\cdot 0{,}004/1{,}004=0{,}00797$.
\end{enumerate}
\resultado \textbf{$P=0{,}4$; $R=0{,}3$; $F_1=0{,}343$. La media armónica castiga al sistema trivial ($0{,}008$) que la aritmética premiaría ($0{,}502$).}
\comprobacion por la \cref{prop:f1}, $F_1=2\cdot 12/(30+40)=24/70=0{,}3429$, entre $R$ y $P$ y más cerca del menor.
\end{solucion}

\begin{solucion}{eje:3-mrr}
\datos $Q_1$: (N, R, N, R, N), $\abs{\Rel}=3$; $Q_2$: (R, R, N, N, N), $\abs{\Rel}=2$.
\planteamiento aplicar la \cref{def:metricas-ranking} con $k=3$.
\pasos
\begin{enumerate}[label=\arabic*.]
  \item $Q_1$: el top-3 contiene 1 relevante; $\Pk{3}=1/3$, $\Rk{3}=1/3$; primer relevante en la posición~2, $\RR=0{,}5$.
  \item $Q_2$: el top-3 contiene 2; $\Pk{3}=2/3$, $\Rk{3}=2/2=1$; $\RR=1$.
  \item Medias: $\Pk{3}=(1/3+2/3)/2=0{,}5$; $\Rk{3}=(1/3+1)/2=0{,}6667$; $\mathrm{MRR}=(0{,}5+1)/2=0{,}75$.
\end{enumerate}
\resultado \textbf{$\Pk{3}$ medio $0{,}5$; $\Rk{3}$ medio $2/3$; $\mathrm{MRR}=0{,}75$.}
\comprobacion $Q_2$, con los relevantes arriba, es mejor que $Q_1$ en las tres métricas, y todos los valores están en $[0,1]$.
\end{solucion}

\begin{solucion}{eje:3-iir-ap}
\datos $\abs{\Rel}=4$; sistema~1: relevantes en las posiciones 1, 3, 9 y 10; sistema~2: en 2, 5, 6 y 7.
\planteamiento $\AP=\frac14\sum_i\Pk{k_i}$; R-precisión $=\Pk{4}$.
\pasos
\begin{enumerate}[label=\arabic*.]
  \item Sistema~1: $\Pk1=1$, $\Pk3=2/3$, $\Pk9=3/9$, $\Pk{10}=4/10$. $\AP_1=(1+0{,}6667+0{,}3333+0{,}4)/4=2{,}4/4=0{,}6$.
  \item Sistema~2: $\Pk2=1/2$, $\Pk5=2/5$, $\Pk6=3/6$, $\Pk7=4/7$. $\AP_2=(0{,}5+0{,}4+0{,}5+0{,}5714)/4=1{,}9714/4=0{,}4929$.
  \item R-precisión: sistema~1, top-4 (R, N, R, N), $2/4=0{,}5$; sistema~2, top-4 (N, R, N, N), $1/4=0{,}25$.
\end{enumerate}
\resultado \textbf{$\AP_1=0{,}6>\AP_2=0{,}493$; R-precisión $0{,}5$ frente a $0{,}25$: ambas prefieren el sistema~1.}
\comprobacion el sistema~1 acierta en la posición~1, y AP pesa mucho los primeros puestos; el sistema~2 recupera los cuatro antes (posición~7 frente a~10), lo que no basta para compensar.
\end{solucion}

\begin{solucion}{eje:3-se}
\datos $\hat p_A=0{,}82$, $\hat p_B=0{,}80$, $n=500$ cada uno, muestras independientes.
\planteamiento la varianza de la diferencia es la suma de varianzas (\cref{prop:se}).
\pasos
\begin{enumerate}[label=\arabic*.]
  \item $0{,}82\cdot 0{,}18/500=0{,}1476/500=0{,}0002952$; $0{,}80\cdot 0{,}20/500=0{,}00032$.
  \item Suma $0{,}0006152$; $\SE=\sqrt{0{,}0006152}=0{,}02480$.
  \item Diferencia $0{,}02$, que equivale a $0{,}02/0{,}0248=0{,}806$ errores estándar.
  \item Para $\SE=0{,}01$: $(0{,}1476+0{,}16)/n=10^{-4}$, luego $n=0{,}3076/10^{-4}=3076$.
\end{enumerate}
\resultado \textbf{$\SE=0{,}0248$; la diferencia es de $0{,}81$ errores estándar, compatible con el azar; hacen falta 3076 preguntas por modelo.}
\comprobacion el error estándar decrece como $1/\sqrt n$: pasar de $0{,}0248$ a $0{,}01$ divide por $2{,}48$ y exige multiplicar $n$ por $2{,}48^{2}=6{,}15$; $500\cdot 6{,}15\approx 3076$.
\end{solucion}

\begin{solucion}{eje:3-iir-ap20}
\datos relevantes en las posiciones 1, 2, 9, 11, 15 y 20 (6 de 20); $\abs{\Rel}=8$; colección de 10\,000.
\planteamiento $\AP=\frac18\sum_i\Pk{k_i}$, con aportación nula de los relevantes no recuperados; en (c), colocar los dos que faltan lo antes y lo más tarde posible.
\pasos
\begin{enumerate}[label=\arabic*.]
  \item (a) $P=6/20=0{,}3$; $R=6/8=0{,}75$; $F_1=2\cdot 0{,}3\cdot 0{,}75/1{,}05=0{,}45/1{,}05=0{,}4286$.
  \item Precisiones en los relevantes: $1$; $1$; $3/9=0{,}3333$; $4/11=0{,}3636$; $5/15=0{,}3333$; $6/20=0{,}3$. Suma: $3{,}3303$.
  \item (b) $\AP=3{,}3303/8=0{,}4163$.
  \item (c) Máxima: posiciones 21 y 22, con $7/21=0{,}3333$ y $8/22=0{,}3636$; $\AP=(3{,}3303+0{,}6970)/8=4{,}0273/8=0{,}5034$.
  \item Mínima: posiciones 9999 y 10\,000, con $7/9999=0{,}0007$ y $8/10\,000=0{,}0008$; $\AP=(3{,}3303+0{,}0015)/8=0{,}4165$.
  \item (d) Error máximo: $0{,}5034-0{,}4163=0{,}0871$.
\end{enumerate}
\resultado \textbf{$P=0{,}3$; $F_1=0{,}429$; $\AP=0{,}416$; rango $[0{,}4165;\,0{,}5034]$; error máximo $\approx 0{,}087$.}
\comprobacion el valor de (b) casi coincide con el mínimo de (c): suponer que lo no recuperado aporta~0 equivale a colocarlo al final del todo.
\end{solucion}

\begin{solucion}{eje:3-passk}
\datos $n$ soluciones independientes, cada una correcta con probabilidad $p$; $c$ correctas; $n\ge k$.
\planteamiento interpretar $\binom{n-c}{k}/\binom{n}{k}$ como la probabilidad de que un subconjunto aleatorio de $k$ soluciones no contenga ninguna correcta.
\pasos
\begin{enumerate}[label=\arabic*.]
  \item Sea $K$ un subconjunto de $k$ de las $n$ soluciones elegido uniformemente al azar e independiente de ellas. Condicionado a que haya $c$ correctas, hay $\binom{n}{k}$ subconjuntos posibles, de los cuales $\binom{n-c}{k}$ contienen solo incorrectas. Luego $P(\text{todas las de }K\text{ incorrectas}\mid c)=\binom{n-c}{k}/\binom{n}{k}$.
  \item Por la ley de la esperanza total, $\E\bigl[\binom{n-c}{k}/\binom{n}{k}\bigr]=P(\text{todas las de }K\text{ incorrectas})$.
  \item Como $K$ es independiente de las soluciones y estas son independientes y equidistribuidas, las $k$ soluciones de $K$ son $k$ ensayos independientes de Bernoulli$(p)$: la probabilidad de que todas fallen es $(1-p)^{k}$.
  \item Por linealidad, $\E\bigl[\widehat{\text{pass@}k}\bigr]=1-(1-p)^{k}$: el estimador es insesgado.
  \item $n=10$, $c=3$: $k=1$, $1-7/10=0{,}3$; $k=2$, $1-21/45=0{,}5333$; $k=5$, $1-\binom75/\binom{10}5=1-21/252=0{,}9167$.
\end{enumerate}
\resultado \textbf{estimador insesgado; con $n=10$ y $c=3$, $0{,}3$, $0{,}5333$ y $0{,}9167$ para $k=1,2,5$.}
\comprobacion con $k=1$ el estimador es $c/n$, la frecuencia de acierto, y crece con $k$, como debe; si $n-c<k$, $\binom{n-c}{k}=0$ y el estimador vale~1, porque cualquier subconjunto de $k$ contiene una correcta.
\end{solucion}

% ---------------------------------------------------------------------
\section{Tema 4}

\begin{solucion}{eje:4-denso}
\planteamiento comparar coincidencia léxica exacta y similitud semántica aprendida (\cref{def:biencoder} y \cref{tema:2}).
\pasos
\begin{enumerate}[label=(\alph*)]
  \item \textbf{Denso.} «automóvil» y «coche» no comparten términos y BM25 no puede relacionarlos; un codificador entrenado con paráfrasis asigna vectores próximos.
  \item \textbf{BM25.} El código «E-404» es un término raro, con idf alta, y la coincidencia exacta lo distingue de «E-403»; un codificador denso puede asignar vectores casi iguales a códigos parecidos.
  \item El índice guarda $E_P(p)$ para toda la colección: si cambia $E_P$, todos esos vectores quedan obsoletos y hay que recodificar la colección. $E_Q$ solo se aplica a la consulta en el momento de buscar. (Ambos deben seguir siendo compatibles, es decir, entrenados para el mismo espacio).
\end{enumerate}
\resultado \textbf{(a) denso; (b) BM25; (c) los vectores de los pasajes están precalculados en el índice.}
\comprobacion es la motivación de la búsqueda híbrida: cada recuperador cubre los fallos típicos del otro.
\end{solucion}

\begin{solucion}{eje:4-coseno}
\datos $q=(2,1,0,1)$; $d_1=(1,1,0,0)$, $d_2=(0,1,1,1)$, $d_3=(2,0,0,1)$.
\planteamiento $\cos(q,d)=q\cdot d/(\norm q\norm d)$ y ordenar de mayor a menor.
\pasos
\begin{enumerate}[label=\arabic*.]
  \item $\norm q=\sqrt6=2{,}449$.
  \item $d_1$: $q\cdot d_1=3$; $\norm{d_1}=\sqrt2$; $\cos=3/\sqrt{12}=0{,}8660$.
  \item $d_2$: $q\cdot d_2=2$; $\norm{d_2}=\sqrt3$; $\cos=2/\sqrt{18}=0{,}4714$.
  \item $d_3$: $q\cdot d_3=5$; $\norm{d_3}=\sqrt5$; $\cos=5/\sqrt{30}=0{,}9129$.
\end{enumerate}
\resultado \textbf{$0{,}9129$ ($d_3$), $0{,}8660$ ($d_1$), $0{,}4714$ ($d_2$); los dos primeros son $d_3$ y $d_1$.}
\comprobacion $q-d_3=(0,1,0,0)$: $d_3$ difiere de $q$ en una sola coordenada; además $\cos(q,d_3)=\sqrt{5/6}=0{,}9129$.
\end{solucion}

\begin{solucion}{eje:4-inbatch}
\datos $S$ del enunciado; positivos en la diagonal.
\planteamiento pérdida de la fila $i$: $-\ln\bigl(e^{S_{ii}}/\sum_je^{S_{ij}}\bigr)$.
\pasos
\begin{enumerate}[label=\arabic*.]
  \item Fila~1, $(3,1,0)$: $e^{3}+e+1=20{,}086+2{,}718+1=23{,}804$; $\pi_1=20{,}086/23{,}804=0{,}8438$; pérdida $0{,}1698$.
  \item Fila~2, $(2,2,1)$: $e^{2}+e^{2}+e=7{,}389+7{,}389+2{,}718=17{,}496$; $\pi_2=7{,}389/17{,}496=0{,}4223$; pérdida $0{,}8620$.
  \item Fila~3, $(0,1,1)$: $1+e+e=6{,}437$; $\pi_3=2{,}718/6{,}437=0{,}4223$; pérdida $0{,}8620$.
  \item Pérdida del lote: $(0{,}1698+0{,}8620+0{,}8620)/3=0{,}6313$.
  \item (b) Cada consulta tiene $B-1=2$ negativos. Las 9 puntuaciones salen de $2B=6$ codificaciones: 3 consultas y 3 pasajes.
  \item (c) Las consultas~2 y~3 empatan: cada una tiene un negativo con la misma puntuación que su positivo ($S_{21}=S_{22}=2$ y $S_{32}=S_{33}=1$), y el modelo no las distingue.
\end{enumerate}
\resultado \textbf{pérdidas $0{,}1698$, $0{,}8620$ y $0{,}8620$; media $0{,}6313$; 2 negativos por consulta y 6 codificaciones; las consultas~2 y~3 son las peor resueltas.}
\comprobacion una pérdida de $\ln 2=0{,}6931$ correspondería a un empate perfecto entre dos candidatos; las filas~2 y~3 superan ese valor porque hay un tercer candidato con probabilidad $0{,}155$.
\end{solucion}

\begin{solucion}{eje:4-pq}
\datos $N=10^{8}$; $d=1024$; $m=64$; $k^\ast=256$; 4 bytes por número.
\planteamiento aplicar la \cref{prop:pq}.
\pasos
\begin{enumerate}[label=\arabic*.]
  \item (a) $10^{8}\cdot 1024\cdot 4=4{,}096\cdot 10^{11}$ bytes $=409{,}6$~GB.
  \item (b) Cada código: $64\log_2 256=512$ bits $=64$ bytes; total $6{,}4\cdot 10^{9}$ bytes $=6{,}4$~GB, 64 veces menos. Libros de códigos: $k^\ast d=262\,144$ números, $1\,048\,576$ bytes $=1$~MiB.
  \item (c) Búsqueda exacta: $Nd=1{,}024\cdot 10^{11}$ multiplicaciones y sumas. Con distancias asimétricas: tablas, $k^\ast d=262\,144$ operaciones; recorrido, $Nm=6{,}4\cdot 10^{9}$ consultas a tabla. Unas 16 veces menos operaciones y 64 veces menos datos leídos de memoria.
  \item (d) $(k^\ast)^{m}=256^{64}=2^{512}\approx 1{,}34\cdot 10^{154}$ reconstrucciones distintas.
\end{enumerate}
\resultado \textbf{(a) $409{,}6$~GB; (b) $6{,}4$~GB de códigos y 1~MiB de libros; (c) $1{,}02\cdot 10^{11}$ frente a $6{,}4\cdot 10^{9}$; (d) $2^{512}$.}
\comprobacion la razón de memoria es $(4\cdot 1024)/64=64$; el número de reconstrucciones es astronómico con solo $64\cdot 256=16\,384$ centroides almacenados.
\end{solucion}

\begin{solucion}{eje:4-hnsw}
\datos $N=10^{6}$; $M=16$; $m_L=1/\ln M$.
\planteamiento aplicar la \cref{prop:hnsw}: $P(\lambda\ge j)=M^{-j}$.
\pasos
\begin{enumerate}[label=\arabic*.]
  \item (a) $m_L=1/\ln 16=1/2{,}7726=0{,}3607$; $P(\lambda\ge 3)=16^{-3}=1/4096=2{,}441\cdot 10^{-4}$.
  \item (b) $NM^{-j}$: capa~1, $62\,500$; capa~2, $3906$; capa~3, $244{,}1$; capa~4, $15{,}26$; capa~5, $0{,}954$.
  \item (c) $\log_{16}10^{6}=6/\log_{10}16=6/1{,}2041=4{,}98$: la capa superior estará en torno a la~5, de modo que hay unas seis capas contando la~0.
\end{enumerate}
\resultado \textbf{(a) $m_L=0{,}3607$ y $P=2{,}44\cdot 10^{-4}$; (b) $62\,500$, $3906$, $244$, $15$ y~$1$; (c) unas seis capas (de la~0 a la~5).}
\comprobacion el máximo nivel esperado exacto es $\sum_{l\ge 1}\bigl(1-(1-16^{-l})^{10^{6}}\bigr)=4{,}68$, coherente con la estimación $\log_M N$.
\end{solucion}

\begin{solucion}{eje:4-iir-dist}
\datos consulta $q$ y documentos $d_i$; distancia euclídea y coseno.
\planteamiento usar el \cref{teo:dist-coseno}.
\pasos
\begin{enumerate}[label=\arabic*.]
  \item (a) Con normas~1, $\norm{q-d_i}^{2}=2-2\cos(q,d_i)$. La función $c\mapsto 2-2c$ es estrictamente decreciente y la raíz cuadrada es creciente, luego $\cos(q,d_i)>\cos(q,d_j)\iff\norm{q-d_i}<\norm{q-d_j}$.
  \item (b) Si solo $\norm{d_i}=1$: $\norm{q-d_i}^{2}=\norm q^{2}+1-2q\cdot d_i$ y $q\cdot d_i=\norm q\cos(q,d_i)$. Como $\norm q$ es la misma para todos los documentos, la distancia es una función estrictamente decreciente del coseno, y los órdenes coinciden.
  \item (c) Contraejemplo: $q=(1,0)$, $d_1=(10,0)$ y $d_2=(0{,}6;\,0{,}8)$. Cosenos: $1$ y $0{,}6$, de modo que el coseno prefiere $d_1$. Distancias al cuadrado: $81$ y $0{,}16+0{,}64=0{,}8$, de modo que la distancia prefiere $d_2$.
\end{enumerate}
\resultado \textbf{con documentos normalizados, distancia creciente y coseno decreciente ordenan igual; sin normalizar, no.}
\comprobacion con $a=(0{,}6;\,0{,}8)$ y $b=(1,0)$ unitarios, $\norm{a-b}^{2}=0{,}8=2-2\cdot 0{,}6$.
\end{solucion}

\begin{solucion}{eje:4-mips}
\datos $M\ge\max_i\norm{x_i}$; $\tilde x_i=(x_i,\sqrt{M^{2}-\norm{x_i}^{2}})$; $\tilde q=(q,0)$.
\planteamiento desarrollar la norma por bloques de coordenadas.
\pasos
\begin{enumerate}[label=\arabic*.]
  \item Como $M^{2}-\norm{x_i}^{2}\ge 0$, la raíz es real y $\tilde x_i$ está bien definido; además $\norm{\tilde x_i}^{2}=M^{2}$ para todo $i$.
  \item $\norm{\tilde q-\tilde x_i}^{2}=\norm{q-x_i}^{2}+\bigl(0-\sqrt{M^{2}-\norm{x_i}^{2}}\bigr)^{2}$.
  \item $=\norm q^{2}-2q^{\top}x_i+\norm{x_i}^{2}+M^{2}-\norm{x_i}^{2}=\norm q^{2}+M^{2}-2q^{\top}x_i$.
  \item $\norm q^{2}+M^{2}$ no depende de $i$; la distancia es una función estrictamente decreciente de $q^{\top}x_i$, y el vecino más próximo es el de máximo producto escalar.
  \item Verificación: $q=(1,1)$, $\norm q^{2}=2$, $M^{2}=9$. Para $x_1=(3,0)$: $\tilde x_1=(3,0,0)$ y $\norm{\tilde q-\tilde x_1}^{2}=4+1+0=5=2+9-6$. Para $x_2=(1,1)$: $\tilde x_2=(1,1,\sqrt7)$ y la distancia es $0+0+7=7=2+9-4$. Para $x_3=(0,2)$: $\tilde x_3=(0,2,\sqrt5)$ y $1+1+5=7$.
\end{enumerate}
\resultado \textbf{$\norm{\tilde q-\tilde x_i}^{2}=\norm q^{2}+M^{2}-2q^{\top}x_i$; MIPS se reduce a vecino más próximo.}
\comprobacion los productos escalares $(3,2,2)$ y las distancias $(5,7,7)$ dan el mismo orden; todos los $\tilde x_i$ tienen norma $M=3$.
\end{solucion}

% ---------------------------------------------------------------------
\section{Tema 5}

\begin{solucion}{eje:5-componentes}
\planteamiento distinguir memoria paramétrica y no paramétrica (\cref{tema:5}).
\pasos
\begin{enumerate}[label=(\alph*)]
  \item \textbf{Falso.} Indexar un documento solo añade su vector al índice; no hay gradiente ni actualización de pesos.
  \item Se entrenan el codificador de consultas y el generador, minimizando $-\ln p(y\mid x)$; el codificador de documentos y el índice se congelan porque actualizarlos exigiría recodificar toda la colección en cada paso.
  \item Cada fragmento recuperado tiene un identificador y una procedencia que pueden citarse; el conocimiento de los pesos está distribuido y no conserva su origen.
\end{enumerate}
\resultado \textbf{(a) falso; (b) codificador de consultas y generador; (c) por la procedencia identificable de los fragmentos.}
\comprobacion si se retira un documento del índice, el sistema deja de poder citarlo sin reentrenar nada.
\end{solucion}

\begin{solucion}{eje:5-troceo}
\datos (a) $L=1000$, $\tau=200$, $s=50$; (b) $L=2300$, $\tau=512$, $s=64$.
\planteamiento inicios $j(\tau-s)$ y número de trozos por la \cref{prop:troceo}.
\pasos
\begin{enumerate}[label=\arabic*.]
  \item (a) Paso $150$; inicios $0,150,300,450,600,750,900$; $n=\lceil 800/150\rceil+1=\lceil 5{,}33\rceil+1=7$. El último cubre $900$--$1099\ge 999$.
  \item (b) Paso $448$; $n=\lceil 1788/448\rceil+1=\lceil 3{,}99\rceil+1=5$; inicios $0,448,896,1344,1792$. El último cubre $1792$--$2303\ge 2299$.
  \item (c) Factores aproximados: $200/150=1{,}33$ y $512/448=1{,}14$. Tokens realmente indexados: $7\cdot 200=1400$ ($1{,}4$ veces) y $5\cdot 512=2560$ ($1{,}11$ veces).
\end{enumerate}
\resultado \textbf{(a) 7 trozos; (b) 5 trozos; (c) unas $1{,}33$ y $1{,}14$ veces más texto.}
\comprobacion con 4 trozos en (b), el último acabaría en $1344+511=1855<2299$: harían falta los 5.
\end{solucion}

\begin{solucion}{eje:5-ragnum}
\datos datos del \cref{ej:rag}.
\planteamiento aplicar la \cref{def:rag}; en RAG-Token, cada factor es una mezcla con pesos $0{,}7$ y $0{,}3$.
\pasos
\begin{enumerate}[label=\arabic*.]
  \item RAG-Sequence: $(a,a)$: $0{,}7\cdot 0{,}8\cdot 0{,}5+0{,}3\cdot 0{,}1\cdot 0{,}2=0{,}286$; $(a,b)$: $0{,}304$; $(b,a)$: $0{,}7\cdot 0{,}2\cdot 0{,}5+0{,}3\cdot 0{,}9\cdot 0{,}2=0{,}124$; $(b,b)$: $0{,}7\cdot 0{,}2\cdot 0{,}5+0{,}3\cdot 0{,}9\cdot 0{,}8=0{,}286$.
  \item Suma: $0{,}286+0{,}304+0{,}124+0{,}286=1$.
  \item RAG-Token: mezcla del primer token, $P(a)=0{,}59$ y $P(b)=0{,}41$; del segundo, $P(a)=0{,}7\cdot 0{,}5+0{,}3\cdot 0{,}2=0{,}41$ y $P(b)=0{,}59$.
  \item $(a,a)=0{,}59\cdot 0{,}41=0{,}2419$; $(a,b)=0{,}3481$; $(b,a)=0{,}41^{2}=0{,}1681$; $(b,b)=0{,}2419$. Suma: $1$.
\end{enumerate}
\resultado \textbf{ambos modelos prefieren $(a,b)$, con $0{,}304$ y $0{,}3481$.}
\comprobacion las marginales del primer token coinciden en los dos modelos ($0{,}286+0{,}304=0{,}59$), como exige la \cref{prop:rag} para $n=1$.
\end{solucion}

\begin{solucion}{eje:5-cmu-f1}
\datos cuatro pares de referencia y predicción; normalización $\nu$ de la \cref{def:em-f1}.
\planteamiento normalizar, comparar cadenas para EM y contar tokens comunes para $F_1$.
\pasos
\begin{enumerate}[label=\arabic*.]
  \item (a) «1900» y «1900»: $\EM=1$; $F_1=1$.
  \item (b) «fundada en 1900» frente a «1900»: $\EM=0$; $c=1$, precisión $1/3$, exhaustividad $1$; $F_1=2\cdot\frac13/\frac43=0{,}5$.
  \item (c) «andrew carnegie y mellon» frente a «andrew carnegie»: $\EM=0$; $c=2$, precisión $2/4$, exhaustividad $1$; $F_1=2\cdot 0{,}5/1{,}5=0{,}6667$.
  \item (d) «profesor graham neubig» (sin el artículo) frente a «graham neubig»: $\EM=0$; $F_1=0{,}8$ (\cref{ej:f1}).
  \item Medias: $\EM=1/4=0{,}25$; $F_1=(1+0{,}5+0{,}6667+0{,}8)/4=0{,}7417$.
\end{enumerate}
\resultado \textbf{$\EM$ medio $0{,}25$ y $F_1$ media $0{,}742$.}
\comprobacion en los cuatro casos la exhaustividad es~1, porque la predicción contiene la referencia; la $F_1$ solo penaliza las palabras sobrantes, y por eso es mayor que $\EM$.
\end{solucion}

\begin{solucion}{eje:5-cs224u}
\datos cinco ejemplos de entrenamiento; criterio doble: exactitud exacta tras normalizar y aparición de la referencia en el contexto.
\planteamiento evaluar las dos condiciones en cada ejemplo.
\pasos
\begin{enumerate}[label=\arabic*.]
  \item (1) «parís» normaliza igual que «París» y el contexto la contiene: \textbf{se conserva}.
  \item (2) La respuesta coincide, pero el contexto no contiene «1492»: \textbf{se descarta}.
  \item (3) «Curie» no coincide exactamente con «Marie Curie»: \textbf{se descarta}.
  \item (4) «El agua» normaliza a «agua» al eliminar el artículo, y el contexto contiene «agua»: \textbf{se conserva}.
  \item (5) «Toronto» no coincide con «Ottawa»: \textbf{se descarta}.
  \item El ejemplo~(2) se acierta con la memoria paramétrica, no con el contexto. Como demostración enseñaría a responder ignorando los pasajes, que es lo contrario de lo que se busca en un sistema RAG.
\end{enumerate}
\resultado \textbf{se conservan los ejemplos~1 y~4 (2 de 5).}
\comprobacion los conservados cumplen las dos condiciones; cada descartado incumple exactamente una.
\end{solucion}

\begin{solucion}{eje:5-descomp}
\datos $a_0=0{,}2$; $k=3$: $r=0{,}7$, $a_1=0{,}9$; $k=10$: $r=0{,}9$, $a_1=0{,}8$.
\planteamiento aplicar la \cref{prop:descomp}.
\pasos
\begin{enumerate}[label=\arabic*.]
  \item (a) $k=3$: $0{,}7\cdot 0{,}9+0{,}3\cdot 0{,}2=0{,}63+0{,}06=0{,}69$. $k=10$: $0{,}9\cdot 0{,}8+0{,}1\cdot 0{,}2=0{,}72+0{,}02=0{,}74$.
  \item (b) $0{,}8=0{,}8r+0{,}2(1-r)=0{,}2+0{,}6r$, luego $r=1$.
  \item (c) Con $a_1=0{,}8$, ni una recuperación perfecta supera $0{,}8$: el generador limita. Conviene mejorar $r$ sin degradar $a_1$, por ejemplo recuperando muchos candidatos y reordenándolos para quedarse con pocos.
\end{enumerate}
\resultado \textbf{(a) $0{,}69$ y $0{,}74$; (b) $r=1$; (c) hay que actuar a la vez sobre la recuperación y sobre la lectura.}
\comprobacion los dos valores de (a) cumplen la cota $r+(1-r)a_0$: $0{,}69\le 0{,}76$ y $0{,}74\le 0{,}92$.
\end{solucion}

\begin{solucion}{eje:5-rag-dem}
\datos \cref{def:rag}; $p_\eta(\cdot\mid x)$ es una distribución sobre $Z_k(x)$ y cada $p_\theta(\cdot\mid x,z,y_{<i})$ lo es sobre $\V$.
\planteamiento sumar sobre $\V^{n}$ intercambiando sumas.
\pasos
\begin{enumerate}[label=\arabic*.]
  \item RAG-Sequence: para cada $z$, $q_z(y)=\prod_i p_\theta(y_i\mid x,z,y_{<i})$ es una distribución sobre $\V^{n}$ por la regla de la cadena (módulo~1). Entonces $\sum_y p_{\mathrm{seq}}(y\mid x)=\sum_z p_\eta(z\mid x)\sum_y q_z(y)=\sum_z p_\eta(z\mid x)=1$, y todos los términos son no negativos.
  \item RAG-Token: el factor $i$, $m_i(y_i\mid y_{<i})=\sum_z p_\eta(z\mid x)p_\theta(y_i\mid x,z,y_{<i})$, es no negativo y $\sum_{y_i}m_i=\sum_z p_\eta(z\mid x)\cdot 1=1$: es una distribución condicional del token $i$ dada la historia. Un producto de condicionales define, por la regla de la cadena, una distribución sobre $\V^{n}$.
  \item Si $n=1$, ambas expresiones son $\sum_z p_\eta(z\mid x)p_\theta(y_1\mid x,z)$.
  \item Si $p_\eta(z^\ast\mid x)=1$, ambas se reducen a $\prod_i p_\theta(y_i\mid x,z^\ast,y_{<i})$.
\end{enumerate}
\resultado \textbf{ambos modelos son distribuciones sobre $\V^{n}$ y coinciden si $n=1$ o si el recuperador es determinista.}
\comprobacion en el \cref{eje:5-ragnum}, con $n=2$ y dos documentos, las sumas valen~1 y los modelos difieren ($0{,}304$ frente a $0{,}3481$ para $(a,b)$).
\end{solucion}

% ---------------------------------------------------------------------
\section{Tema 6}

\begin{solucion}{eje:6-vf}
\planteamiento contrastar con la \cref{def:herramienta}, la documentación de Claude y la especificación de MCP.
\pasos
\begin{enumerate}[label=(\alph*)]
  \item \textbf{Falsa.} El modelo solo emite texto; el sistema valida y ejecuta.
  \item \textbf{Verdadera.} El bloque \texttt{tool\_result} lleva el \texttt{tool\_use\_id} del bloque \texttt{tool\_use}~\cite{claude-tools}.
  \item \textbf{Falsa.} MCP es un protocolo, basado en JSON-RPC~2.0, entre aplicaciones y servidores de datos y herramientas~\cite{mcp-spec}.
  \item \textbf{Verdadera}, entre sus principios de seguridad~\cite{mcp-spec}.
\end{enumerate}
\resultado \textbf{(a) F; (b) V; (c) F; (d) V.}
\comprobacion (a) y (d) expresan la misma idea: lo que llega al modelo, incluidas las descripciones de herramientas, es texto que puede ser malicioso, y quien actúa es el sistema.
\end{solucion}

\begin{solucion}{eje:6-correo}
\datos agente que resume correos; un correo contiene una orden de reenvío a una dirección externa.
\planteamiento la orden llega dentro de un dato que el agente lee (\cref{def:inyeccion}).
\pasos
\begin{enumerate}[label=\arabic*.]
  \item Es una \textbf{inyección de instrucciones indirecta}: el atacante escribe en una fuente que el agente consume.
  \item Defensa~1, mínimo privilegio: el resumidor solo tiene permiso de lectura y ninguna herramienta de reenvío. Límite: si otra tarea del mismo agente necesita reenviar, el riesgo reaparece.
  \item Defensa~2, confirmación humana antes de enviar nada fuera del dominio. Límite: la fatiga; tras muchas confirmaciones se acepta sin leer.
  \item Complemento: delimitar cada correo como dato y avisar de que no contiene órdenes. Ayuda, pero el modelo puede ignorarlo.
\end{enumerate}
\resultado \textbf{inyección indirecta; mínimo privilegio y confirmación humana, ambas parciales.}
\comprobacion con permiso solo de lectura, aunque el modelo obedezca la orden, ningún correo sale del buzón.
\end{solucion}

\begin{solucion}{eje:6-diseno}
\datos tres capacidades: leer correos, crear eventos y enviar mensajes.
\planteamiento asignar permisos según el daño posible y tratar todo contenido externo como no fiable.
\pasos
\begin{enumerate}[label=\arabic*.]
  \item Leer correos: libre, en solo lectura. Es la fuente del riesgo: cualquier remitente puede escribir instrucciones.
  \item Crear eventos: como borradores o en un calendario aparte; el daño es reversible y basta con un aviso.
  \item Enviar mensajes: acción irreversible y vía de exfiltración; exige confirmación humana mostrando el texto y el destinatario, y una lista blanca de destinatarios.
  \item Separación: tras leer contenido externo, el agente no puede enviar en el mismo turno sin pasar por el usuario. Se registran todas las llamadas para auditarlas.
  \item El \emph{prompt} no garantiza nada: instrucciones y datos llegan en la misma secuencia de tokens y el modelo genera la continuación más probable; un texto persuasivo puede cambiarla, y ninguna frase del \emph{prompt} lo impide con certeza.
\end{enumerate}
\resultado \textbf{lectura libre, eventos como borrador y envíos con confirmación y lista blanca: la seguridad está en el sistema.}
\comprobacion en el peor caso, con el modelo completamente engañado, solo puede proponer un mensaje que el usuario ve y puede rechazar: el daño queda acotado.
\end{solucion}

\begin{solucion}{eje:6-fiab}
\datos éxito por paso $p$, independiente.
\planteamiento por la \cref{prop:fiabilidad}, la probabilidad de completar $n$ pasos es $p^{n}$; se despeja con logaritmos.
\pasos
\begin{enumerate}[label=\arabic*.]
  \item (a) $0{,}95^{2}=0{,}9025$; $0{,}95^{4}=0{,}8145$; $0{,}95^{8}=0{,}6634$; $0{,}95^{10}=0{,}6634\cdot 0{,}9025=0{,}5987$.
  \item (b) $p^{20}=0{,}8\Rightarrow\ln p=\ln 0{,}8/20=-0{,}2231/20=-0{,}01116\Rightarrow p=0{,}9889$.
  \item (c) $0{,}95^{n}<0{,}5\iff n>\ln 0{,}5/\ln 0{,}95=13{,}51$ (al dividir por un número negativo se invierte la desigualdad): $n=14$.
  \item (d) Con un reintento, el éxito por paso es $1-0{,}05^{2}=0{,}9975$, y en diez pasos $0{,}9975^{10}=0{,}9753$.
\end{enumerate}
\resultado \textbf{(a) $0{,}599$; (b) $p\ge 0{,}9889$; (c) a partir de 14 pasos; (d) $0{,}975$.}
\comprobacion $0{,}95^{13}=0{,}5133>0{,}5>0{,}95^{14}=0{,}4877$; y en (b), $1-20\cdot 0{,}0111=0{,}778\approx 0{,}8$, la aproximación lineal para fallos pequeños.
\end{solucion}

\begin{solucion}{eje:6-creencia}
\datos $b(\text{sí})=0{,}8889$ tras la primera búsqueda; $P(\text{enc.}\mid\text{sí})=0{,}8$, $P(\text{enc.}\mid\text{no})=0{,}1$; el estado no cambia.
\planteamiento aplicar la \cref{prop:creencia} con $T$ igual a la identidad.
\pasos
\begin{enumerate}[label=\arabic*.]
  \item (a) «encontrado»: $b''(\text{sí})=0{,}8\cdot 0{,}8889/(0{,}8\cdot 0{,}8889+0{,}1\cdot 0{,}1111)=0{,}7111/0{,}7222=0{,}9846$.
  \item (b) «no encontrado», con verosimilitudes $0{,}2$ y $0{,}9$: $b''(\text{sí})=0{,}2\cdot 0{,}8889/(0{,}2\cdot 0{,}8889+0{,}9\cdot 0{,}1111)=0{,}1778/0{,}2778=0{,}64$.
  \item (c) Como la búsqueda no modifica el estado, $T(s'\mid s,a)=\mathbb{1}[s'=s]$ y $\sum_sT(s'\mid s,a)b(s)=b(s')$: la actualización es la regla de Bayes con la creencia anterior como distribución a priori.
\end{enumerate}
\resultado \textbf{(a) $0{,}9846$; (b) $0{,}64$; (c) el estado es estático.}
\comprobacion (b) es exacto: $0{,}5\cdot 0{,}8\cdot 0{,}2/(0{,}5\cdot 0{,}8\cdot 0{,}2+0{,}5\cdot 0{,}1\cdot 0{,}9)=0{,}08/0{,}125=0{,}64$, partiendo de la creencia inicial y aplicando las dos observaciones a la vez.
\end{solucion}

\begin{solucion}{eje:6-toolformer}
\datos probabilidades de los tres tokens en los tres casos; $\tilde w_t=\max(0,1-0{,}2t)$.
\planteamiento normalizar los pesos, calcular las tres pérdidas y aplicar la \cref{def:toolformer}.
\pasos
\begin{enumerate}[label=\arabic*.]
  \item (a) $\tilde w=(1;\,0{,}8;\,0{,}6;\,0{,}4;\,0{,}2;\,0;\dots)$, de suma~3. Luego $w_0=0{,}3333$, $w_1=0{,}2667$ y $w_2=0{,}2$.
  \item Logaritmos: $-\ln 0{,}10=2{,}3026$; $-\ln 0{,}12=2{,}1203$; $-\ln 0{,}90=0{,}1054$; $-\ln 0{,}30=1{,}2040$; $-\ln 0{,}60=0{,}5108$; $-\ln 0{,}50=0{,}6931$.
  \item Sin llamada: $0{,}76753+0{,}32106+0{,}13863=1{,}2272$.
  \item Llamada sin resultado: $0{,}70676+0{,}32106+0{,}13863=1{,}1664$.
  \item (b) Con resultado: $L_i^{+}=0{,}03512+0{,}13622+0{,}13863=0{,}3100$.
  \item $L_i^{-}=\min(1{,}2272;\,1{,}1664)=1{,}1664$, y la diferencia es $L_i^{-}-L_i^{+}=0{,}8565$.
  \item (c) Con $\tau_f=1$, $0{,}8565<1$: se descarta. Con $\tau_f=0{,}5$: se conserva.
\end{enumerate}
\resultado \textbf{$w=(0{,}3333;\,0{,}2667;\,0{,}2)$; $L_i^{+}=0{,}3100$, $L_i^{-}=1{,}1664$, diferencia $0{,}8565$; se descarta con $\tau_f=1$ y se conserva con $\tau_f=0{,}5$.}
\comprobacion el tercer token tiene la misma probabilidad en los tres casos y aporta lo mismo, $0{,}1386$; la diferencia procede de los dos primeros, los más próximos a la llamada.
\end{solucion}

\begin{solucion}{eje:6-horizonte}
\datos $H=3$; estados sin información (S) e informado (I); responder acierta con $0{,}5$ en S y $0{,}9$ en I; buscar pasa de S a I con probabilidad $0{,}5$.
\planteamiento inducción hacia atrás sobre los pasos restantes $h$, con $V_h(s)$ el valor óptimo; después, evaluar las cuatro políticas estacionarias.
\pasos
\begin{enumerate}[label=\arabic*.]
  \item $h=1$ (paso~3): buscar no deja tiempo para responder y vale~0. $V_1(\mathrm S)=0{,}5$ y $V_1(\mathrm I)=0{,}9$: responder.
  \item $h=2$ (paso~2): en I, responder vale $0{,}9$ y buscar, $V_1(\mathrm I)=0{,}9$; $V_2(\mathrm I)=0{,}9$. En S, responder vale $0{,}5$ y buscar, $0{,}5\,V_1(\mathrm I)+0{,}5\,V_1(\mathrm S)=0{,}45+0{,}25=0{,}7$; $V_2(\mathrm S)=0{,}7$: buscar.
  \item $h=3$ (paso~1): en S, buscar vale $0{,}5\cdot 0{,}9+0{,}5\cdot 0{,}7=0{,}8>0{,}5$; $V_3(\mathrm S)=0{,}8$.
  \item (a) Política óptima: buscar en S en los pasos~1 y~2, responder en S en el paso~3 y responder en I. Valor $0{,}8$.
  \item (b) Políticas estacionarias: si $\pi(\mathrm S)=$ responder, el valor es $0{,}5$. Si $\pi(\mathrm S)=\pi(\mathrm I)=$ buscar, nunca responde y vale~0. Si $\pi(\mathrm S)=$ buscar y $\pi(\mathrm I)=$ responder: con probabilidad $0{,}5$ se informa en el paso~1 y acierta con $0{,}9$ en el~2; si no, con probabilidad $0{,}5$ se informa en el paso~2 y acierta con $0{,}9$ en el~3; si no, busca en el paso~3 y obtiene~0. Valor: $0{,}5\cdot 0{,}9+0{,}25\cdot 0{,}9=0{,}675$.
  \item El máximo entre las estacionarias es $0{,}675<0{,}8$: la acción óptima en S depende del paso (buscar en los pasos~1 y~2, responder en el~3).
  \item (c) Un agente con un límite de pasos debe cambiar de conducta al acercarse al límite: dejar de explorar y responder con lo que tenga. La información del paso debe estar en su contexto; con ella, la política deja de ser estacionaria en el estado físico.
\end{enumerate}
\resultado \textbf{valor óptimo $0{,}8$, con una política que depende del paso; la mejor política estacionaria obtiene $0{,}675$.}
\comprobacion si el horizonte fuera infinito y sin descuento, buscar siempre en S acabaría informando al agente con probabilidad~1 y la política estacionaria «buscar en S, responder en I» alcanzaría $0{,}9$: la dependencia del paso procede del horizonte finito.
\end{solucion}
